\section{Introduction}
Three dimensional (3D) object reconstruction and acquisition are a long time focus of machine vision research and a key to visual understanding and interpretation. It has a wide range of practical applications, including robotics, AR/VR, autonomous navigation, and industrial automation, to name just a few.
Being able to access the 3D geometric cue about an object or scene can immensely benefit in post-processing and decision making for instance in recognition \cite{pontil1998support}, segmentation \cite{qi2017pointnet}, pose estimation \cite{andriluka2010monocular} and other computer vision and pattern recognition tasks.

As a computer vision task, this requires a sensor for observing the object, which can be of two types, namely, active and passive cameras.
Active cameras like Laser imaging detection and ranging (LiDAR), Time of Flight (ToF), and depth cameras like the Kinect directly acquire scene depth and, in turn, the 3D shape of an object but come with their limitations apart from the high cost and power requirements.
Passive cameras are much cheaper and use significantly less power to acquire two-dimensional (2D) optical RGB snapshots of the object, which can then be used to recover the 3D structure by posing this as an inverse problem.

Recovering 3D shape from its two-dimensional snapshots has been a goal of classic computer vision problems like multiview stereo and Shape-from-X techniques.
Early methods approached this with a geometric perspective and focused on mathematically modeling the process of 2D projection to pose it as an ill-posed problem \cite{hartley2003multiple} constrained by suitable scene priors.
The structure from motion(SfM) techniques can be used to obtain a sparse 3D point cloud of an object along with the respective camera poses.
Along similar lines, Multi-View Stereo (MVS) techniques can provide a dense 3D reconstruction.
Recently, with the advent of deep learning-based methods, a model can directly learn a non-linear mapping from the images to the 3D representation outperforming classical methods while doing it in real-time \cite{pix2vox}.

% MVS to NBV
Although there are works that attempt to reconstruct the 3D shape of an object from just a single view but these are in early stages due to their highly ill-posed nature of the problem.
Moreover, as in MVS, taking more than one view of the object reduces ambiguity and occlusion to significantly improve the 3D recognition quality as more information is available for recovery.
Typically, these views are acquired from a set of passively selected viewpoints independent of the object geometry.
To obtain the best 3D reconstruction, one would expect these viewpoints to depend on the 3D shape in consideration.
Therefore, to obtain the best possible 3D reconstruction quality from the captured images, one must actively decide the next best viewpoint to take the snapshot from.

% Motivation and Intro to NBV
The Next View Selection (NVS) problem has been explored since the early days of machine vision research \cite{connolly1985determination}.
This problem arises in many scenarios where one would want to find the next best viewpoint around an object or in an arena while maximizing a particular quality or performance metric.
For instance, to better grasp an object, a robotic arm equipped with a camera may want to find the next viewpoint to obtain the best 3D reconstruction for better understanding quickly.
Similarly, a UAV hovering over a scene may want to find the next view, which can provide the best 3D digital elevation model to minimize the flight time.
Several works in robotic navigation, robotic arm positioning, and remote sensing are discussed in the related work section.

% Gaps in the literature, typically done and what they lack
However, almost all of these works use range cameras to capture scene depth directly and select the next view either heuristically or by optimizing for the surface area covered, or information gained.
Along with this, they have aimed to reconstruct a surface representation of the scene like mesh or a point cloud.
Volumetric representation like voxels is better suited for practical applications due to ease of manipulation and training than surface or mesh-based representations.
Hence, we focus on learning a volumetric representation of the scene in this work.
Moreover, selecting the next views using only passive camera images of a scene has been explored much recently but not very actively in the literature. We focus on this particular setup in this paper.
Predicting a 3D shape or next views using passive cameras is a more challenging task due to much lesser data available for use than in active sensors.
Due to this setup, classical solutions based on the surface reconstruction and optimized for the surface area covered or information gained cannot be used here.

% A hint of advantages of the proposed method and how it fits into the bigger picture
This paper attempts to address these gaps and propose a novel classification-based approach for NVS without using the preprocessed ground truth next views.
Since our goal is to maximize the 3D reconstruction quality, we propose to extract a supervisory signal from the reconstruction process itself to drive the classification model's training and obtain a better loss-goal correspondence.
We present how a light-weight neural network-based classifier can be incorporated into existing 3D reconstruction models to perform NVS and trained without the knowledge of ground truth NVS labels.
Training such a classifier-reconstructor architecture is not straightforward without the labels. Therefore, we propose a novel loss objective that facilitates end to end training of such a joint classifier-reconstructor deep neural network model.

% End with a hint of some key findings and insights gained
We demonstrate how the proposed method successfully selects the next views to maximize the reconstruction quality for a given 3D object on synthetic and real data.
Additionally, we analyze the selected next views for each object category to gain insights and observe underlying patterns to learn how they are distributed depending on the object shape.
Our contributions in this paper are as follows:
\begin{itemize}
    \item We propose a deep learning-based approach for next view selection and show how it can be effectively approached as a classification problem.
    \item We show how the supervisory signal from the 3D reconstruction process can be used to learn both the classification and reconstruction networks jointly in an end to end manner.
    \item We present a qualitative analysis of the selected next views to gain insights into the process and their dependence on object categories and shape.
    \item We present exhaustive experiments on both synthetic and real data to demonstrate the effectiveness of the proposed method in selecting appropriate next views while providing high reconstruction quality.
\end{itemize}

\begin{figure*}
\begin{center}
\includegraphics[scale=0.36]{figures/3d-nvs-overview.pdf}
\end{center}
   \caption{An Overview of the proposed 3D-NVS model used during the (a) training and (b) testing. A detailed explanation of each sub-parts is given in Section \ref{sec:method}. (c) Chosen viewpoints of the ShapeNet dataset to generate images for training.}
\label{fig:overview-model}
\end{figure*}
